\begin{problemsummary}
{A client that has not changed starts getting \code{Cannot query field} or
\code{argument is required}, at HTTP 200, with an \code{errors} array and no
\code{data} key, after somebody else deployed a subgraph. Or nothing visible
happens at all: a field that was non-null is nullable now, every response is
still correct, and a null arrives months later in a client with no branch for
it.}
{Composition reads the subgraph documents and nothing else, so every question
it can answer is a question about them. Whether a field is still being
selected is not one: the operation lives in a mobile release, a dashboard or a
script, outside the repository the documents live in. Of fourteen one-line
edits composed against this book's graph, twelve pass at exit zero with
nothing printed, and every one of those twelve is on the reference
implementation's breaking-change list. The two the composer refuses are both
routes it could not plan, and one of them breaks no client at all.}
{Diff the composed schema rather than the subgraph document: compose what is
published, compose the candidate, and diff \code{engineConfig.graphqlSchema}
out of the two configs, which costs two commands you already run and reports
the change in the schema clients actually read. Classify that change using
the compatibility rules; a text diff does not do the classification. Retain
old client operations and replay them against the candidate before deployment.
That catches breaks in the retained corpus even without a control plane, but
cannot cover operations you never collected. Observed traffic supplies another
corpus; both vendors' default seven-day window still leaves older or rarely
used operations outside the evidence.}
\end{problemsummary}
